\subsection{二阶张量的推广}
\label{sec:app-a-second-order-tensors}
\renewcommand{\thestatement}{A.2.\arabic{statement}}
\renewcommand{\theHstatement}{appA.sec02.\arabic{statement}}
\setcounter{statement}{0}

我们的目标并不是对张量作详细的一般性研究. 本书只需要详细考虑二阶张量场, 即 $\mathbb R^d$ 的自同态场.

因此, 我们借助 $\mathbb R^d$ 的典范基, 直接把二阶张量等同于方阵.

\begin{Definition}
\label{def:app-a-tensor-gradient-divergence-laplacian}
设 $\Omega$ 是 $\mathbb R^d$ 的一个区域.
\begin{itemize}
 \item 设 $v$ 是 $\Omega$ 上的正则向量场; 称如下张量场为 $v$ 的梯度:
 \[
  \nabla v=\left(\frac{\partial v_i}{\partial x_j}\right)_{1\leq i,j\leq d}.
 \]
 \item 设 $\sigma=(\sigma_{ij})_{1\leq i,j\leq d}$ 是 $\Omega$ 上的正则张量场; 称如下向量场为 $\sigma$ 的散度:
 \[
  \operatorname{div}\sigma
  =\left(\sum_{j=1}^{3}\frac{\partial\sigma_{ij}}{\partial x_j}\right)_{1\leq i\leq d}.
 \]
 \item 设 $v$ 是 $\Omega$ 上的正则向量场; 称如下向量场为 $v$ 的 Laplace 算子:
 \[
  \Delta v=\operatorname{div}(\nabla v).
 \]
\end{itemize}
容易验证, $\Delta v$ 是这样的向量场: 在 $\mathbb R^d$ 的一个固定标准正交基中, 它的各分量为 $\Delta v_i$, 其中 $v_i$ 是向量场 $v$ 在同一基中的各分量.
\end{Definition}

\begin{Definition}{Tensor product}
\label{def:app-a-tensor-product}
设 $u,v$ 是两个向量场; 称如下定义的二阶张量为 $u$ 和 $v$ 的张量积, 并将它记为 $u\otimes v$: 在一个标准正交基中,
\[
 u\otimes v=(u_iv_j)_{1\leq i,j\leq d}.
\]
若把 $u\otimes v$ 作用于第三个向量场 $w$, 则有
\[
 (u\otimes v)\mathbin{.}w=(v\cdot w)u.
\]
此外, 张量积的散度有如下公式:
\begin{equation}
 \operatorname{div}(u\otimes v)=(\operatorname{div}v)u+(v\cdot\nabla)u.
 \label{eq:app-a-divergence-tensor-product}
\end{equation}
\end{Definition}

最后, 可以在二阶张量场的集合上定义一个标量积, 它就是把矩阵空间等同于 $\mathbb R^{d\times d}$ 后该空间中的 Euclidean 标量积. 在张量语言中, 这称为缩并张量积.

\begin{Definition}{Contracted product}
\label{def:app-a-contracted-product}
设 $\sigma,\tau$ 是两个二阶张量; 将它们的缩并积定义为如下实数:
\[
 \sigma:\tau=\sum_{1\leq i,j\leq d}\sigma_{ij}\tau_{ij}.
\]
与这个标量积相联系的范数直接记为 $|\cdot|$, 从而
\[
 |\sigma|^2=\sigma:\sigma.
\]
\end{Definition}
